\chapter{How the Virtual Machine Works}
\label{ch:vm}

\standfirst{Enough of the implementation to explain why the system behaves as
it does. You do not need this to write Smalltalk; you need it the first time
something surprises you.}

\section{The shape}

\begin{plain}
src/port/       threads, mutexes, condition variables, TLS, time, atomics
src/om/         object memory (bb and mt), image readers, collector
src/interp/     bytecode interpreter, contexts, method lookup, primitives
src/gfx/        BitBlt, display, SDL3 pump, the rasterised face
src/sched/      Process, Semaphore, the scheduler
src/compiler/   Smalltalk compiler in C, chunk and Tonel readers
src/boot/       image bootstrap
\end{plain}

About 40,000 lines of C11 with no dependency but SDL3, and that only for the
window.

\section{Objects}

Every object is a header plus fields, and every reference to one goes through
an \textbf{object table}: an oop is an index, not an address.

Small integers are the exception---they are tagged directly in the oop, so
arithmetic on them allocates nothing.

\subsection{Why keep an object table?}

Spur and every other modern VM dropped it, and under threading the trade
inverts:

\begin{itemize}
\item \ct{become:}---swapping two objects' identities---is one atomic swap of
  two table entries, instead of a stop-the-world scan of the heap for every
  reference.
\item Pinning an object for SDL or a foreign call is free: the body does not
  move because nothing points at it directly.
\item Compaction touches one entry per object rather than every reference to
  it.
\end{itemize}

The cost is one indirection per access. The purchase is atomic identity
mutation, which is what a parallel Smalltalk needs.

\subsection{Two object memories}

\ct{src/om/om.h} defines the object-memory operations as macros, and exactly
one implementation is compiled in, so the hot path inlines either way.

\ct|om_bb.c| is Blue Book Chapter 30 verbatim: 16-bit oops, reference
counting. It exists so the \emph{same interpreter} can load the real 1983
image and reproduce Xerox's traces. \ct|om_mt.c| is the real one: 64-bit
table, tagged integers, per-worker allocation, a marking collector at a
safepoint, weak references and ephemerons.

Passing \ct{trace2} therefore validates the same code the parallel system
runs. That is the central trick of the project.

\section{Contexts are objects}

A stack frame is a \ct{MethodContext}: an ordinary heap object with a sender,
an instruction pointer, a stack pointer, a method, a receiver and slots. There
is no machine stack.

That single decision is why so much is possible: processes are cheap because a
process is a pointer to a context chain; the debugger can show you the stack
because the stack is data; you can edit a method in the middle of a running
call and restart the frame; and \ct{ensure:} can find its own marker by
walking senders.

\section{Bytecodes}

Squeak's V3PlusClosures set. About 250 opcodes, most of them one byte,
covering push, store, pop, jump, send, and return.

Two things worth knowing as a user:

\textbf{Special selectors.} \ct{+}, \ct{-}, \ct{<}, \ct{=}, \ct{at:},
\ct{at:put:}, \ct{do:} and about two dozen others have a bytecode of their own.
They are still real sends---the fast path checks the receiver's class and
falls through to a full lookup if it is not what it expected---so overriding
\ct{+} in your class works.

\textbf{Inlined control flow.} \ct{ifTrue:}, \ct{and:}, \ct{whileTrue:} and
\ct{to:do:} with literal blocks compile to jumps. No block is created and no
message is sent, which is why they cost nothing and why a non-Boolean receiver
reports \ct{must be a boolean} from a jump rather than a doesNotUnderstand.

\section{Primitives}

A method may declare \ct{<primitive: N>}. The VM tries N; if it succeeds the
Smalltalk body never runs, and if it fails the body runs as a fallback. That
fallback is why \ct{SmallInteger>>+} can hand off to \ct{LargePositiveInteger}
arithmetic without the VM knowing anything about large integers.

Two primitives in this system are not operations at all but \emph{labels}: 198
and 199 mark unwind and handler frames, and nothing implements them, so they
always fail and the Smalltalk body runs. Signalling finds them by walking the
sender chain. That is how exceptions are written in Smalltalk.

Primitives 240--255 are reserved for this system's parallel operations.

\section{Method lookup}

Receiver's class, then its method dictionary, then up the superclass chain to
\ct{Object}, whose superclass is \ct{nil}. If nothing implements the selector,
\ct{doesNotUnderstand:} is sent instead.

\begin{stnote}[There is no method cache]
Primitive 89, \ct{flushCache:}, is a no-op because there is nothing to flush.
That is a measured decision rather than an omission: \ct|lookup_method| probes
from the selector's hash, and its inclusive cost is 0.1\% of a steady-state
profile and 8.4\% of a whole image build---the most send-heavy thing the system
does. A \emph{perfect} cache would buy 8.4\% there and nothing elsewhere, at
the price of a per-worker cache, a global epoch, and method dictionaries made
immutable-once-published, whose failure mode is a stale method silently
answering an old body. Reference counting is 25\% of the same profile. That is
where the next measurement goes.
\end{stnote}

\section{The collector}

Marking, at a safepoint. The interpreter polls for safepoints at message sends
and backward jumps; when one is requested, every worker parks, one thread
collects, and everybody resumes. No Smalltalk code observes a partially
collected heap.

Reference counts are rebuilt from the root walk, which has a consequence that
bit the bootstrap for months: \textbf{a reference held only in C protects
nothing}. Anything the VM holds must be visited by the root walk, or the
collector will reclaim it and hand its table slot to the next allocation.

The shape of that bug is worth recognising. A message sent to an object of a
class that has no business being there---\ct|aMethodContext does not understand #priority|---is rarely a wrong send. It is almost always a freed
object whose table slot has been handed to something else.

Three rules follow from it that the server's gates enforced. A primitive
that makes more than one object makes the container first and pushes it on
the stack before making the elements, because an allocation can run a
collection and the stack is a root where a C variable is not. A slot that
every worker writes---the scheduler's \ct{activeProcess}---is exchanged
atomically, because a plain counted store lets two writers release the same
old value. And \ct{someInstance} and \ct{nextInstance}, which walk the
object table, do so at a safepoint when workers are running: a freed header
is reused in place, so a walk cannot read one while another worker refills
it. That makes \ct{allInstances} stop the world once per instance under
\ct{-serve}; it is a tool, not a path.

\section{The bootstrap}

There is no chicken-and-egg problem, because the compiler is in C. Building an
image is: read every class definition, create the classes, compile every
method, run the class initialisers, construct the handful of objects that
exist only in a live image---\ct{Sensor}, \ct{ScheduledControllers},
\ct{Transcript}, \ct{Smalltalk}, the top \ct{Project}, \ct{Disk}---and write
the object memory out.

That last group is the interesting part, and it is why a bootstrapped image is
not merely a compiled library: those objects exist in the 1983 image because
somebody made them interactively in 1979 and never stopped saving, and a
bootstrap has to construct each one deliberately.

\section{The window}

Thread 0 pumps SDL and never executes Smalltalk. Workers execute Smalltalk and
never call SDL video. SDL3's main-thread-only rules and macOS's Cocoa run loop
force that shape, and it also avoids a real deadlock: a thread parked in a GC
safepoint at the moment the window server wants a response would hang the
compositor.

The display Form is the authoritative pixel store; the window is a view of it.
Thread 0 copies damaged regions out and presents them, and never writes back.

\section{Where to read next}

\begin{center}
\small
\begin{tabular}{@{}ll@{}}
\toprule
\ctp{doc/CONCURRENCY.md} & the contract, the primitive table, the locking rules\\
\ctp{doc/SCALING.md} & the benchmark and what each kernel measures\\
\ctp{doc/Display.md} & the window, the face, antialiasing\\
\ctp{doc/LanguageExtensions.md} & every post-1983 syntax and where it stands\\
\ctp{doc/PLAN-TO-PHARO.md} & where the project is going\\
\ctp{doc/LICENSING.md} & what may be redistributed\\
\bottomrule
\end{tabular}
\end{center}
